\chapter{Physical and Mathematical Framework}
\label{ch:theory}
\section{Geometry and Nondimensionalization}
The default S-channel joins two semicircular segments with a continuous
tangent and piecewise curvature. Its nominal radii are
$R_{\rm in}=\InnerRadius$ and $R_{\rm out}=\OuterRadius$; the centerline
radius is their average. The width is $W=R_{\rm out}-R_{\rm in}=\ChannelWidth$
and the centerline length is approximately $L=\ChannelLength$.
A general supported channel is represented by
\begin{equation}
 \vect{x}(s,\eta)=\vect{c}(s)+h\eta\vect{n}(s),\qquad
 0\le s\le L,\quad -1\le\eta\le1,\quad h=W/2,
 \label{eq:geometry}
\end{equation}
where $\vect{c}$ is a planar centerline and $\vect{n}$ its unit normal.
Curvature and overlap checks exclude invalid offset tubes. A custom CSV
changes the centerline, not the physical topology supported by the solver.

Dimensional variables are scaled using $L_0$, $U_0$ and $\mu U_0/L_0$ for
length, speed and pressure. Time is scaled by $T_0=L_0/U_0$. Throughout the
following equations, variables without a dimensional subscript are
dimensionless. The reference Reynolds number is
$\Rey=\rho_f U_0 L_0/\mu$; $L_0$ is not implicitly a hydraulic diameter.
The current principal case uses $\Rey=\ReynoldsNumber$.

\ResultFigure{01_geometry}{Channel and sampling}{Default channel and sampled
interior and boundary points from the repository. Coordinates are
nondimensional; the displayed point set illustrates the sampling procedure.}{fig:geometry}

\section{Fluid Equations and Boundary Conditions}
For a steady incompressible Newtonian fluid,
\begin{align}
 \nabla\cdot\vect{u}&=0,\label{eq:continuity}\\
 \Rey(\vect{u}\cdot\nabla)\vect{u}&=-\nabla p+\nabla^2\vect{u}.
 \label{eq:momentum}
\end{align}
The walls satisfy no slip. Both end sections have the prescribed tangential
velocity $\vect{u}=(1-\eta^2)\vect{t}$ with unit center speed. Pressure has
one gauge condition. In particular, the outlet is a prescribed-velocity
boundary, rather than an assumed zero-pressure outlet. The FEM and PINN use
these same physical boundary conditions, with their pressure constants
aligned during comparison. The prescribed section flux is $Q_0=4h/3$.

\section{Particle and Fluid-Marker Dynamics}
Let $d_p$ be particle diameter and $\rho_r=\rho_p/\rho_f$ the density ratio.
The adopted drag-only heavy-particle equations are
\begin{align}
 \frac{d\vect{x}_p}{dt}&=\vect{v}_p,\label{eq:particle-position}\\
 \frac{d\vect{v}_p}{dt}&=\frac{18}{\rho_r\Rey d_p^2}
 [\vect{u}(\vect{x}_p)-\vect{v}_p]
 =\frac{\vect{u}(\vect{x}_p)-\vect{v}_p}{\St},
 \qquad \St=\frac{\rho_r\Rey d_p^2}{18}.
 \label{eq:particle-velocity}
\end{align}
The particle Reynolds number is
$\Rey_p=\Rey d_p\lVert\vect{u}-\vect{v}_p\rVert$.
Stokes drag requires $\Rey_p\ll1$ and a particle small relative to flow
length scales. The force formulation is a simplification of small-sphere
dynamics, not the full Maxey--Riley equation \cite{maxey1983}.

A fluid marker follows $d\vect{x}/dt=\vect{u}(\vect{x})$. A finite-size
particle sticks when its center reaches a distance $d_p/2$ from the wall;
its velocity is then set to zero. This contact rule is irreversible and
does not represent rebound or adhesion-force calibration. Initial particle
velocity equals the local fluid velocity unless the user supplies a
different admissible value.

\section{Force-Network Interpretation}
The normalized drag force $\vect{F}^{*}$ is scaled by displaced-fluid mass
times $U_0^2/L_0$. With input
$z=(u,v,p,d_p,x_p,y_p,u_p,v_p)$ and normalized input map $N$, the implemented
surrogate is
\begin{equation}
 \widehat{\vect{F}}^{*}(z)=\frac{18}{\Rey d_p^2}
 \{g_\phi(Nz)-g_\phi(Nz_0)\},
 \label{eq:force-network}
\end{equation}
where $z_0$ replaces the particle velocity by the fluid velocity. Subtraction
enforces zero force at zero slip. The network is trained against the
algebraic Stokes residual, so it approximates an already specified law.
Pressure and particle position are input slots; they do not establish
additional force mechanisms. Exact Stokes drag remains the reference.

\section{Physical Limits}
The model omits gravity, lift, Brownian motion, added mass, history effects,
fluid feedback, particle collisions, and near-wall lubrication corrections.
The largest tested diameter is 15\% of the channel width, so that case is a
sensitivity study rather than proof of the unbounded point-particle model's
physical accuracy. The computed maximum $\Rey_p$ is
\MaximumParticleRe; this supports a low-slip-inertia regime but does not
remove confinement or dilute-loading assumptions.
